\chapter{Ingredients for adaptive corruptions}
\label{chap:corruption}

Section~\ref{sec:corruptions} describes what a simulator must explain when the prover is corrupted
during or after the protocol of \cite{JKO13}. This chapter formalizes the protocol-specific facts
behind that description: receiver equivocality of the PVW OT in decryption mode
\cite[Sect.~3.3, Theorem~3.2]{GKPS18} and of the Chou--Orlandi OT for a party that knows the
sender's secret; the explanation of the prover's OT messages and of its revealed label for any
witness \cite[Sect.~3.2]{GKPS18}; the equivocation of the hash commitment by programming the
random oracle; and soundness with the PVW OT in decryption mode, the mode in which it is receiver
equivocal. None of these statements is new. Adaptive security itself, which needs a model of
execution with corruptions such as the one of universal composability \cite[Sect.~7.1]{Can20}, is
not defined here.

\section{Receiver-equivocal oblivious transfer}

\begin{definition}[Receiver-equivocal OT]
  \label{def:re-ot}
  \lean{GcZk.AC.ReceiverEquivocal}
  \leanok
  Let the receiver of an OT send $\mathit{msg}(\sigma, r)$ for the choice bit $\sigma$ and the
  randomness $r$. It is \emph{receiver equivocal}, with simulated message $\mathit{sim}$ and
  explanations $e_0, e_1$, if each $e_\sigma$ is a bijection of the randomness and
  \[ \mathit{msg}(\sigma, e_\sigma(r)) = \mathit{sim}(r) \quad \text{for all } \sigma, r . \]
\end{definition}

A simulator sends $\mathit{sim}(r)$ for a uniform $r$, without the choice bit; told $\sigma$ later,
it explains its message by the randomness $e_\sigma(r)$, which is uniform because $e_\sigma$ is a
bijection. This is \cite[Sect.~3, Definition~3.1]{GKPS18}, where the explained randomness is only
required to be statistically close to the honest one; here it has exactly the same distribution.

\begin{theorem}[The PVW receiver in decryption mode]
  \label{thm:pvw-re}
  \lean{GcZk.AC.pvw_dec_equivocal, GcZk.AC.pvw_dec_equivocal'}
  \leanok
  \uses{def:re-ot, def:pvw-dec}
  On a decryption-mode common string, with $g_1 = y \cdot g_0$, $y \neq 0$ and $h_b = x \cdot g_b$,
  the PVW receiver is receiver equivocal with simulated message $\mathit{keyGen}(0, r)$,
  $e_0(r) = r$ and $e_1(r) = y^{-1} r$.
\end{theorem}

\begin{proof}
  \leanok
  By Definition~\ref{def:pvw-dec}, $\mathit{keyGen}(1, r) = \mathit{keyGen}(0, y r)$, so
  $\mathit{keyGen}(1, y^{-1} r) = \mathit{keyGen}(0, r)$; multiplication by $y^{-1}$ is a bijection
  of $\mathbb{Z}_q$. This is the proof of \cite[Theorem~3.2]{GKPS18} and the algorithm
  $\mathrm{TrapKeyGen}$ of \cite[Sect.~5.3]{PVW08}.
\end{proof}

\begin{proposition}[The Chou--Orlandi receiver, given the sender's secret]
  \label{prop:co-re}
  \lean{GcZk.AC.co_equivocal}
  \leanok
  \uses{def:re-ot, def:co-ot}
  Let $A = a \cdot g$. The receiver message $B = b \cdot g + wA$ is receiver equivocal with
  simulated message $b \cdot g$ and explanations $e_0(b) = b$, $e_1(b) = b - a$.
\end{proposition}

\begin{proof}
  \leanok
  $(b - a) \cdot g + a \cdot g = b \cdot g$, and $b \mapsto b - a$ is a bijection.
\end{proof}

The explanation needs $a$, which the simulator knows when it plays the honest verifier, and
otherwise learns only at the open-all step (Section~\ref{sec:corruptions}).

\section{Explaining the prover}

\begin{theorem}[Explaining the prover]
  \label{thm:explain-prover}
  \lean{GcZk.AC.explainRand, GcZk.AC.explainRand_bijective, GcZk.AC.explain_prover}
  \leanok
  \uses{def:re-ot, thm:zk}
  Let the OT receiver be receiver equivocal, and let the verifier choose its message (the garbled
  circuit, the OT inputs and the seed it opens) as a function of the prover's OT messages. Let
  $C(w) = 1$, and let $r_j$ be the simulator's randomness for the $j$-th OT. With the explained
  randomness $e_{w_j}(r_j)$ the prover's OT messages are the simulated ones $\mathit{sim}(r_j)$,
  and the verifier's view of the prover is the output of the simulator of Theorem~\ref{thm:zk}.
  The explained randomness is uniform when $r$ is.
\end{theorem}

\begin{proof}
  \leanok
  The first claim is the definition, coordinate by coordinate. Then the verifier's messages are the
  same in both cases, and Theorem~\ref{thm:zk} gives the second claim. The explanation is a
  bijection in each coordinate.
\end{proof}

With an ideal commitment this is all a simulator must explain when the prover is corrupted after
the execution \cite[Sect.~3.2]{GKPS18}: every step of the prover after the oblivious transfers is
independent of the witness.

\section{Equivocating the hash commitment}

With the hash commitment of Chapter~\ref{chap:commitment}, a prover corrupted after committing and
before opening hands over its output label $Z_w$ and its commitment randomness $r$. In the real
execution its commitment is $G(Z_w, r)$. The simulator commits instead to a uniform $c$, without the
witness, and at the corruption reveals $(Z_w, r)$ and sets $G(Z_w, r) := c$; this is possible in
the \emph{programmable} random oracle model \cite[Sect.~1.2]{GKPS18}. We reuse the games of
Definition~\ref{def:com-games}, with a check that always passes, so that the last message of the
real game is $(Z_w, r)$.

\begin{theorem}[Equivocation by programming the oracle]
  \label{thm:equivocate-rom}
  \lean{GcZk.AC.always, GcZk.AC.simC, GcZk.AC.hyb_eq_simC, GcZk.AC.card_realC_le,
    GcZk.AC.card_simC_le, GcZk.AC.equivocate_rom}
  \leanok
  \uses{def:com-games, lem:hyb-real, lem:ans-update, thm:zk-rom}
  Let a malicious verifier make $q_1$, $q_2$ and $q_3$ oracle queries before the commitment,
  between the commitment and the corruption, and after the corruption. Over $G$, $c$ and $r$
  uniform, the probabilities that it outputs $1$ in the real execution and against the simulator
  differ by at most $(q_1 + q_2 + q_3)/|R|$.
\end{theorem}

\begin{proof}
  \leanok
  Off the bad event $\mathrm{Bad}_{23}$ of Lemma~\ref{lem:hyb-sim}, the hybrid of
  Lemma~\ref{lem:hyb-real}, which programs the oracle already before the second phase, and the
  simulator differ only in the answers of the second phase, and these do not hit the programmed
  point (Lemma~\ref{lem:ans-update}). The rest is the counting of Theorem~\ref{thm:zk-rom}.
\end{proof}

\section{Soundness in decryption mode}

The PVW OT is receiver equivocal in decryption mode, while the soundness of
Theorem~\ref{thm:pvw-sound} uses messy mode, where the sender is statistically protected. Switching
the modes costs the DDH advantage of the acceptance probability.

\begin{theorem}[Soundness with the PVW OT in decryption mode]
  \label{thm:pvw-dec-sound}
  \lean{GcZk.AC.crsM_messy, GcZk.AC.pAcc, GcZk.AC.pAcc_nonneg, GcZk.AC.pAcc_le_one,
    GcZk.AC.prob_accept_dec_le}
  \leanok
  \uses{thm:pvw-sound, thm:pvw-modes}
  Let $C$ be unsatisfiable, $q > 2$, and let the cheating prover depend on the common string and
  make at most $N_1$ and $N_2$ queries in its two phases. Averaged over the decryption-mode setup,
  the verifier accepts with probability at most
  \[ \frac{N_1 + 1}{|V| - 1} + \frac{N_1^2 + N_2 + 1}{|D|} + \mathrm{Adv}_{\mathrm{DDH}}(p_{\mathrm{acc}})
     + \frac{3}{q} , \]
  where $p_{\mathrm{acc}}$ is the acceptance probability as a function of the common string.
\end{theorem}

\begin{proof}
  \leanok
  The strings of the messy setup are messy, so by Theorem~\ref{thm:pvw-sound} the acceptance
  probability on each of them is at most the first two terms, and so is its average. The
  acceptance probability lies in $[0, 1]$, and Theorem~\ref{thm:pvw-modes} bounds the difference of
  the two averages.
\end{proof}
